\section{The graph estimate}\label{graph:section}

The independent angular energy uses a common radial frequency and two
independent directions.  We now pass from a pair of depth-$a$ cells to
depth-$p$ cells.  After splitting the endpoint transforms into cell
pieces, angular and radial oscillation restrict the four cell labels to
a graph of small row and column degrees.  The parallel tests give the
degree bound.

The cell terms remain oscillatory until the nonstationary terms have
been removed.  We then take absolute values of the retained terms and
bound them by positive independent angular energies on the smaller
cells.  At the end of the Section, Cauchy--Schwarz transfers this
estimate to the separated form $\mathsf H_F$ with the same endpoint
symbols.

\begin{proposition}[Graph estimate for the independent angular energy]
\label{graph:estimate}
Let $\rho_1,\rho_2$ be two regularized class laws, possibly equal, and
write $M_j$, $A_{j,n}$, $I_j$ and $k_j$ for their profiles and associated
quantities from Section~\ref{test:section}.  Let $X,Y$ be occupied
cells of depth $a$.  Put $X_1=X$, $X_2=Y$, and let
$b^{(j)}_z(w)$ be an admissible endpoint symbol on $X_j$ in the sense
of Definition~\ref{osc:symbols}.  Let $L$ be the largest nominal
length of a descriptor in either symbol, with $L=0$ if both lists are
empty.  Suppose that the integer depths satisfy $0\le a<p\le v\le N$
and
\begin{equation}\label{graph:range}
 p-a\le v-p+2\sigma N,\qquad
 L\le v-p+2\sigma N,\qquad v-p\ge10\sigma N.
\end{equation}
At endpoint $j$, suppose that the symbol imposes the parallel test
$[g_j,p]$ at a level $\ell_j\in\{0,\ldots,J_0\}$: for
$\rho_j$-almost every $z\in X_j$ and every $w\in S^{d-1}$,
\begin{equation}\label{graph:support-test}
 b^{(j)}_z(w)\ne0
 \quad\Longrightarrow\quad
 w\in G^{(j)}_{\ell_j,[g_j,p]}(z).
\end{equation}
The superscript $j$ indicates the class used in the test.  Assume
\begin{equation}\label{graph:anchor-range}
 a\le g_j<p,\qquad g_j\le a+2\sigma N.
\end{equation}

Form the general independent energy $\mathsf H^{X,Y}(a,v)$ of
Definition~\ref{osc:states} using these specified symbols.  For every
fixed $T>0$,
\begin{equation}\label{graph:bound}
 \begin{split}
 \mathsf H^{X,Y}(a,v)
 &\le
 R^{-\sum_{j=1}^2(A_{j,p}-A_{j,a})
       +\min_{j=1,2}k_j(a,p)+O_d(\sigma)+o(1)}\\
 &\qquad\cdot
 \max_{\substack{P\subset X,\ Q\subset Y\\
                   P,Q\text{ occupied depth-}p\text{ cells}}}
       \mathsf H^{P,Q}(p,v)
 +O(R^{-T}).
 \end{split}
\end{equation}
No separation of $X$ and $Y$ is required.  Every child in
\eqref{graph:bound} is formed by restricting the same endpoint symbols
to $P$ and $Q$.  It retains every original mask factor, including any
whose anchor is $p$.  Constants are uniform over cells and allowed
symbols with fixed list-size exponents and derivative budgets, and with
fixed common numerical bounds for their smooth angular factors at the
derivative orders used.  The $o(1)$ tends to zero as $N\to\infty$ with
these bounds, the parameters, and $T$ fixed.
\end{proposition}

We give the localization and degree arguments separately.  Set
\begin{equation}\label{graph:scales}
 h=p-a,\qquad q=v-p,\qquad
 \delta=\min(c_0,2^{-h}R^{10\sigma}),\qquad
 \tau=2^{-p}R^{30\sigma},
\end{equation}
where $c_0>0$ is a sufficiently small dimensional constant.  Take a
nonnegative smooth partition of unity $\{\eta_\nu\}$ on the sphere,
with bounded overlap, such that $\eta_\nu$ is supported in a cap of
radius $C\delta$ centered at $w_\nu$, and
\begin{equation}\label{graph:bin-derivatives}
 \|\nabla^m\eta_\nu\|_\infty\le C_m\delta^{-m}.
\end{equation}
This can be obtained by normalizing scaled bumps on a finite atlas;
the number of caps is $O_d(\delta^{-(d-1)})$.  No relation between the
two cap centers used below is imposed.

For occupied depth-$p$ cells, use the transforms $U_P,V_Q$ of
Equation~\eqref{osc:transforms}, divided by the model masses $M_1(p)$
and $M_2(p)$, respectively.  The exact splitting identities are
\begin{equation}\label{graph:splitting}
 U_X=\lambda_1\sum_{P\subset X}U_P,\qquad
 V_Y=\lambda_2\sum_{Q\subset Y}V_Q,
 \qquad \lambda_j=\frac{M_j(p)}{M_j(a)}=R^{-I_j(a,p)}.
\end{equation}
Expand the two squares in $\mathsf H^{X,Y}(a,v)$ and insert the angular
partition in both variables.  Index each term by a row $(P,Q)$ and a
column $(P',Q')$, with
$x\in P$, $x'\in P'$, $y\in Q$, and $y'\in Q'$.  The $P$-labels belong
to $\rho_1$ and the $Q$-labels to $\rho_2$.  Put
\[
 \Delta_1=x-x',\qquad \Delta_2=y-y',\qquad
 d_1=c_P-c_{P'},\qquad d_2=c_Q-c_{Q'}.
\]
With the spatial points fixed, the term has phase
\begin{equation}\label{graph:phase}
 r\big(w_1\cdot\Delta_1+w_2\cdot\Delta_2\big).
\end{equation}
Each displacement joins two points in the same depth-$a$ parent cell.
Consequently
\begin{equation}\label{graph:displacement-size}
 |\Delta_j|+|d_j|\lesssim_d2^{-a}.
\end{equation}
The locations of $X$ and $Y$ relative to one another do not enter this
bound.

For an endpoint cell $P$ with measure $\rho$ and symbol $b_x$, define
\[
 E(P,\nu)=\int_P\int_{S^{d-1}}
                  \eta_\nu(w)|b_x(w)|^2\,dw\,d\rho(x).
\]
Here symbols are set to zero on any fixed spatial null set where their
stated support conditions have not been specified.  The cell--cap
incidence $(P,\nu)$ is \emph{eligible} precisely when $E(P,\nu)>0$.
If $E(P,\nu)=0$, Fubini's Theorem shows that, for almost every $w$ with
$\eta_\nu(w)>0$, the symbol vanishes for $\rho$-almost every $x\in P$.
The corresponding transform then vanishes for every $r$, so all
expanded terms involving this incidence have zero integral and may be
omitted.  Positive $E(P,\nu)$ supplies a point--angle witness where the
symbol is nonzero and the stated support conditions hold.  At each
fixed $N$ only finitely many witnesses are chosen; no measurable
selection or lower bound for their masses is needed.

Fix two caps, with centers $w_1^0,w_2^0$.  Retain an edge precisely when
all its cell--cap incidences are eligible and
\begin{equation}\label{graph:retained-tests}
 \big|\operatorname{proj}_{(w_j^0)^\perp}d_j\big|\le\tau
 \quad(j=1,2),\qquad
 |w_1^0\cdot d_1+w_2^0\cdot d_2|\le\tau.
\end{equation}
Eligibility and the edge decision concern an entire cell--cap term.
They never cut its angular support by a nonsmooth pointwise condition.

\begin{lemma}[Nonstationary localization]\label{graph:localization}
Under \eqref{graph:range}, replacing the expanded independent energy by its
sum over the retained edges changes it by $O(R^{-T})$, for every fixed
$T>0$.  This replacement is made in the integrated oscillatory form,
before taking absolute values of the expanded cell terms.
\end{lemma}

\begin{proof}
Fix the four spatial points in one cell term.  Fubini's Theorem is
applicable, since the spatial measures are finite and all amplitudes
are bounded with compact radial support.  Replacing centers by points,
or replacing cap centers by angles in their cap supports, changes
each test in \eqref{graph:retained-tests} by at most
\begin{equation}\label{graph:test-error}
 E\lesssim_d2^{-p}+2^{-a}\delta
 \lesssim_d2^{-p}R^{10\sigma}=o(\tau).
\end{equation}
The second inequality also holds when $\delta=c_0$, since that case
implies $c_0\le2^{-h}R^{10\sigma}$.  Fixed geometric constants are
absorbed by the power gap between $R^{10\sigma}$ and $R^{30\sigma}$.

Suppose first that the last inequality in
\eqref{graph:retained-tests} fails.  Throughout this entire term,
\[
 |w_1\cdot\Delta_1+w_2\cdot\Delta_2|\ge\tau/2
\]
for sufficiently large $N$.  In the variable $u=r/2^v$, each radial
integration by parts gains a factor at most $C(2^v\tau)^{-1}$.
The symbols and angular partition do not depend on $u$, so only the
fixed smooth radial cutoff is differentiated.  Since
\begin{equation}\label{graph:oscillation-size}
 2^v\tau=2^qR^{30\sigma}\ge R^{40\sigma},
\end{equation}
repeating this operation gives arbitrary power decay.

Otherwise, one of the two perpendicular inequalities fails.  Suppose
it is the first.  Let $t$ be the unit tangent vector at $w_1^0$ in the
direction of $\operatorname{proj}_{(w_1^0)^\perp}d_1$, and extend it
by the smooth sphere vector field
\[
 V(w)=t-(t\cdot w)w.
\]
All derivatives of this vector field are bounded independently of the
cap.  By \eqref{graph:test-error},
\[
 |V(w_1)\cdot\Delta_1|\ge\tau/2
\]
throughout its support.  Applied to the phase
$\Phi=r(w_1\cdot\Delta_1+w_2\cdot\Delta_2)$, the field therefore
satisfies
\[
 |V\Phi|\gtrsim2^v\tau,
 \qquad |V^m\Phi|\lesssim_m2^v2^{-a}\quad(m\ge1).
\]
The same reasoning applies in the second variable if its perpendicular
test fails.

Here are the derivative bounds needed to iterate this integration by
parts.  By Definition~\ref{osc:symbols}, Lemma~\ref{test:product-symbols}
and \eqref{graph:bin-derivatives}, the product amplitude has angular
derivatives of order $m$ bounded by
$C_m(1+D_*+\delta^{-1})^m$, where
\[
 D_*=R^\sigma\max(1,2^LR^{-50\sigma}).
\]
Differentiating $(V\Phi)^{-1}$ contributes, in addition to its initial
factor $(2^v\tau)^{-1}$, at most powers of
$1+2^{-a}/\tau$.  Indeed, every positive-order derivative of $V\Phi$
is bounded by $C_m2^v2^{-a}$, whereas $|V\Phi|\gtrsim2^v\tau$.
The bounded divergence and derivatives of $V$ cause only constant
factors.  Thus $m$ integrations by parts give a bound by a constant
$C_m$ times
\begin{equation}\label{graph:ibp-ratio}
 \left(
 \frac{1+D_*+\delta^{-1}+2^{-a}/\tau}{2^v\tau}
 \right)^m
\end{equation}
for the angular--radial integral of the fixed spatial term.
This follows directly by induction with the formal adjoint of
$(iV\Phi)^{-1}V$; a smooth extension of the reciprocal on a slightly
larger cap may be used, since its denominator has the same lower
bound there.

The three hypotheses in \eqref{graph:range} imply
\begin{align*}
 D_*&\le R^\sigma+2^qR^{-47\sigma},\\
 \delta^{-1}&\le c_0^{-1}+2^qR^{-8\sigma},\\
 2^{-a}/\tau&=2^hR^{-30\sigma}\le2^qR^{-28\sigma}.
\end{align*}
Consequently the numerator in \eqref{graph:ibp-ratio} is at most
$C2^qR^{4\sigma}$, and the ratio is at most $CR^{-26\sigma}$.
There is therefore a fixed positive power gain per integration by
parts.  Taking sufficiently many derivatives gives $O(R^{-T'})$
for any prescribed $T'$.

Finally, all sums and normalizations outside the fixed spatial term
cost at most a fixed power of $R$.  There are at most $R^{4d}$ cell
quadruples, at most $O(R^{2(d-1)})$ cap pairs, the profiles lie between
$0$ and $d$, and the energy prefactors are at most $R^d$.  Spatial
integrations contribute at most the same kind of fixed power through
their normalizations.  Choose $T'$ larger than $T$ by these fixed
powers before summing.  The required derivative order depends on the
fixed parameters and $T$, but is independent of $N$.  No spatial
derivative has been taken, so bounded measurable spatial multipliers
are allowed.  This proves the assertion.
\end{proof}

\begin{lemma}[Both degrees of the retained graph]\label{graph:degrees}
For each fixed cap pair, both the row degrees and the column degrees
of the retained graph are at most
\begin{equation}\label{graph:degree-bound}
 D_{\mathrm{gr}}=
 R^{\min_{j=1,2}k_j(a,p)+O_d(\sigma)+o(1)}.
\end{equation}
This estimate does not require the two cap directions to be separated.
\end{lemma}

\begin{proof}
Fix a row $(P,Q)$.  Its unknown column members $P'$ and $Q'$ have
profiles $1$ and $2$, respectively.  Consider the unknown of profile
$j$.  Its perpendicular constraint places its center within $\tau$ of
a line through its fixed mate in direction $w_j^0$.  Group its eligible
candidate cells by their depth-$g_j$ ancestor.  Since the candidate
cells lie in one depth-$a$ parent, the number of these groups is at
most
\[
 2^{d(g_j-a)}\le R^{2d\sigma}.
\]
In one nonempty group choose an eligible cell $P_*$ and a point--angle
witness $(z_*,w_*)$ from its incidence.  The support condition
\eqref{graph:support-test} gives its profile's parallel test
$[g_j,p]$ at $z_*\in P_*$.  This test is based at the center $c_*$ of
the anchor cell $P_*$.  In applications where a
normalized derivative of a mask occurs, the same implication holds by
Lemma~\ref{test:smooth-masks}.

For every other candidate center $c$ in this group, the two
perpendicular constraints and $|w_*-w_j^0|\lesssim\delta$ give
\begin{equation}\label{graph:witness-tube}
 |c-c_*|\le\sqrt d\,2^{-g_j},\qquad
 \operatorname{dist}(c,c_*+\mathbb Rw_*)
 \lesssim_d2\tau+2^{-g_j}\delta
 \lesssim_d2^{-p}R^{30\sigma}.
\end{equation}
For large $N$, the latter is smaller than every test width
$W_i(p)\ge2^{-p}R^{500\sigma}$.  The first inequality lies inside the
testing ball $B(c_*,10(d+1)2^{-g_j})$.  The parallel test therefore
bounds the number of candidates in this group by
\[
 R^{k_j(g_j,p)+C_d\sigma}
 \le R^{k_j(a,p)+C_d\sigma}.
\]
The last inequality holds because $[g_j,p]\subseteq[a,p]$.
Summing the groups bounds the number of possibilities for this first
unknown by $R^{k_j(a,p)+O_d(\sigma)}$.

Fix such a first unknown.  Its displacement, and hence its contribution
to the scalar constraint in \eqref{graph:retained-tests}, is now fixed.
That scalar constraint confines the longitudinal coordinate of the
remaining unknown to an interval of length $O(\tau)$ in its own cap
direction.  Its perpendicular constraint confines its other $d-1$
coordinates to a ball of radius $O(\tau)$.  Together they confine its
center to a ball of radius $O_d(\tau)$, containing at most
\[
 C_d(1+2^p\tau)^d\lesssim_d R^{30d\sigma}
\]
depth-$p$ grid centers.  We may begin with either profile $j$.
Taking the better of the resulting two counts proves the asserted
row-degree bound.

For a fixed column, the unknown row members again have one fixed mate
each, and the signs of their displacements are reversed.  Group either
unknown by its own depth-$g_j$ ancestors and repeat
\eqref{graph:witness-tube}.  Once that member is fixed, the scalar and
perpendicular constraints confine the other to the same $O(\tau)$ ball.
This gives the identical column-degree bound.  The scalar constraint
always fixes the remaining coordinate in that unknown's own cap
direction, so no division by an angle between $w_1^0$ and $w_2^0$
occurs.
\end{proof}

\begin{proof}[Proof of Proposition~\ref{graph:estimate}]
Apply Lemma~\ref{graph:localization}.  For fixed $r,w_1,w_2$ in a cap
pair, take absolute values only of the retained terms.  In the
row--column indexing above, the amplitudes are
\[
 F_{P,Q}=U_P(r,w_1)V_Q(r,w_2),\qquad
 G_{P',Q'}=U_{P'}(r,w_1)V_{Q'}(r,w_2).
\]
The elementary inequality $2|FG|\le|F|^2+|G|^2$, together with both
degree bounds in Lemma~\ref{graph:degrees}, gives
\begin{equation}\label{graph:schur}
 \sum_{(\alpha,\beta)\text{ edge}}|F_\alpha G_\beta|
 \le\frac{D_{\mathrm{gr}}}{2}
       \left(\sum_\alpha|F_\alpha|^2+
             \sum_\beta|G_\beta|^2\right).
\end{equation}
The sums on the right may be enlarged to all row and column labels.
Sum the nonnegative partition weights over cap pairs.  Their sums
are one, so there is no factor equal to the number of caps.  Integrate
against the same nonnegative radial cutoff $\Psi(r/2^v)\,dr/2^v$.
Each row square gives the independent-angle energy of its child cell
pair, and each column square does so as well.  Thus
\eqref{graph:schur} involves only child quantities
$\mathsf H^{P,Q}(p,v)$ formed from the restricted symbols.

Let $I_j=I_j(a,p)$.  The occupied-cell estimate in
Lemma~\ref{test:profile-bounds} bounds the number of row labels, and
also of column labels, by
\[
 R^{I_1+I_2+O_d(\sigma)+o(1)}.
\]
The splitting factors in \eqref{graph:splitting} appear squared,
giving $R^{-2(I_1+I_2)}$.  Finally, the parent prefactor divided by
the child prefactor is
\[
 \frac{2^{d(v-a)}}{2^{d(v-p)}}=2^{d(p-a)}.
\]
Since $S=d/2$ and
$I_j=A_{j,p}-A_{j,a}+S(p-a)/N$, their product is
\begin{align*}
 R^{-2(I_1+I_2)}R^{I_1+I_2}2^{d(p-a)}
 &=R^{-(I_1+I_2)+d(p-a)/N}\\
 &=R^{-\sum_{j=1}^2(A_{j,p}-A_{j,a})}.
\end{align*}
Multiplying by \eqref{graph:degree-bound} proves
\eqref{graph:bound}, including the arbitrary-power error from
Lemma~\ref{graph:localization}.  The argument has only restricted
symbols to smaller cells, so every child retains all its original
mask factors and support conditions.
\end{proof}

\begin{corollary}[The separated angular form]\label{graph:separated-form}
Under the hypotheses of Proposition~\ref{graph:estimate}, suppose also
that $X,Y$ and their supported laws have the separated geometry used
to define $\mathsf H_F^{X,Y}(a,v)$ in Section~\ref{osc:section}.  Then
\eqref{graph:bound} holds with its left-hand side replaced by
$\mathsf H_F^{X,Y}(a,v)$ and with the same child energies.  This includes
the outer root $a=0$ with separated actual supports inside possibly
equal depth-zero cells $X,Y$.
\end{corollary}

\begin{proof}
Equation~\eqref{osc:angular-cs} gives
$\mathsf H_F^{X,Y}(a,v)\le\mathsf H^{X,Y}(a,v)$ with the same endpoint
symbols.  Their tests and model-mass normalizations are therefore
unchanged.  Proposition~\ref{graph:estimate} supplies exactly the child
energies in \eqref{graph:bound}.
\end{proof}
